\begin{abstract}
\noindent
Eight Glacial Lake Outburst Floods (GLOFs) from the Plaine Morte glacier (2011--2018)
affected the Simme catchment before a supraglacial bypass channel ended the outburst
regime in 2019. We test whether these events distort flood frequency statistics using
two long-term discharge records at the Simme gauge in Lenk: daily maximum discharge
(BAFU/FOEN, 1944--2025) and daily mean discharge (GRDC, 1944--2020). Under daily
maxima, GLOFs produced the annual maximum in 4 of 8 GLOF years; under daily means in
6 of 8, because sustained outflows elevate the 24-hour average above brief storm peaks.
The 2018 outburst ($\hat{T} \approx 789$\,yr) dominates the statistical signal.
Kolmogorov--Smirnov and Rank-Sum tests reject homogeneity between the GLOF-inclusive
record and the pre-GLOF reference in both datasets ($p \leq 0.002$). After removing
direct GLOF contributions with a $\pm 5$-day masking window, both tests fail to reject
homogeneity ($p \geq 0.17$), indicating the natural-flood regime remained stationary.
A consistent but non-significant upward tendency in the masked peaks suggests possible
gradual hydroclimatic change that warrants monitoring as post-bypass data accumulate.
Naive inclusion of GLOF peaks overstates the 100-year design flood by approximately
15\,\%.
\end{abstract}
